
Our experiments demonstrate that task heterogeneity in multi-domain NLP benchmarks creates a 15.09\% relative improvement oracle gap (11.62 percentage points absolute) between per-task optimal adapter rank selection and the best fixed-rank baseline. This section interprets our findings and acknowledges limitations.

\subsubsection{Key Findings and Interpretation}

The measured oracle gap of 15.09\% relative improvement (11.62 percentage points absolute) provides quantitative evidence that no single fixed adapter configuration can serve all tasks optimally in multi-domain deployment scenarios.

First, it validates the assumption underlying task-aware adapter routing research: the optimization opportunity is substantial (>15\% relative improvement upper bound), not a marginal edge case. Before investing in complex routing mechanisms, evidence is needed that the gap exists and is worth exploiting. Our results provide that evidence.

Second, it challenges current practice in parameter-efficient fine-tuning. Practitioners typically choose a single LoRA rank (commonly 8 or 16) based on validation performance on one task or average performance across a task suite, then apply it uniformly. Our results show this approach leaves substantial performance on the table when deploying across heterogeneous task distributions.

Third, it establishes an upper bound for task-aware routing mechanisms. The 15.09\% relative improvement oracle gap assumes perfect hindsight selection with zero errors. Practical routing mechanisms will have classifier errors, regret from errors, and overhead costs. If a meta-learned routing policy can recover even 60\% of the oracle gap while managing these challenges, it would represent substantial value. Our measurement provides the benchmark against which future routing approaches should be evaluated.

The even distribution of oracle selections across ranks (5/4/4/4) is a surprising finding. If most tasks preferred rank-8 with a few outliers selecting other ranks, the oracle gap might reflect noise. The uniform distribution proves that different task types genuinely require different capacity levels.

This finding extends prior work on task diversity (Wang et al., 2018; Hu et al., 2020) from linguistic phenomena to adapter capacity requirements. GLUE and XTREME were designed to span diverse domains, dataset sizes, and linguistic phenomena—our results show that this diversity manifests in heterogeneous optimal configurations.

The systematic patterns we observe (Chinese tasks prefer rank-4, German tasks prefer rank-32, small datasets prefer low ranks) suggest that task meta-features provide signal for adapter selection. This is encouraging for future routing mechanisms: if optimal rank correlates with observable task characteristics, lightweight classifiers could learn these patterns.

Rank-32's collapse to random baseline (50\%) on datasets up to 67K samples is a stark demonstration of overfitting. The critical finding is that this collapse occurs not just on tiny datasets (WNLI: 635 samples) but also on medium-sized datasets like SST-2 with 67,349 samples. With 262,144 adapter parameters and 67K samples, the capacity-data ratio is approximately 1:256—one parameter per 256 training samples. This should theoretically be sufficient, yet rank-32 still collapses, suggesting that LoRA's capacity requirements are more stringent than full fine-tuning.

This finding has practical implications for rank selection guidelines. Literature consensus (Hu et al., 2021) suggests ranks 4-16 work well for most tasks, but rarely documents what happens at higher ranks with medium-sized datasets. Our systematic evaluation shows that rank >16 not only provides diminishing returns but can actively harm performance through overfitting when dataset size is below 300K samples.

However, we must acknowledge ambiguity about whether rank-32's poor performance reflects fundamental overfitting or hyperparameter mismatch from our uniform training protocol. With rank-specific regularization (higher dropout, stronger weight decay, more aggressive early stopping), rank-32 might perform better. Resolving this requires future experiments with rank-specific hyperparameter tuning.

\subsubsection{Limitations}

Our work has several limitations that bound the scope of our claims.

\textbf{Limitation 1: Oracle Gap vs Routing Benefit}

We measure oracle gap (15.09\% relative improvement, 11.62 percentage points absolute) under perfect hindsight selection, not realistic routing benefit achievable by imperfect classifiers. Oracle gap represents an upper bound, not achievable improvement. Practical routing mechanisms will have classifier errors, regret from errors, and overhead costs. With 70\% routing accuracy, net benefit ≈ (oracle gain × accuracy) - (regret from errors) - (overhead) ≈ 6-8\% absolute improvement, not the full 15.09\% relative improvement oracle gap.

This is an existence hypothesis, designed explicitly as foundation validation before investing in routing mechanisms. The oracle gap proves the optimization opportunity exists and is substantial—mechanism development is the next phase, not a limitation of the current work.

\textbf{Limitation 2: Single-Seed Directional Validation}

All experiments use single random seed (42) without confidence intervals or statistical significance testing. We cannot claim statistical significance, only directional evidence. The oracle gap magnitude (15.09\% relative improvement) and systematic patterns (uniform oracle distribution, language-specific correlations) suggest robustness despite single-seed validation.

Multi-seed validation with 95\% confidence intervals is computationally expensive (68 configurations × N seeds) and deferred to mechanism validation.

\textbf{Limitation 3: Accuracy-Based Oracle (Not Multi-Objective)}

Oracle selection uses accuracy only, ignoring FLOPs, latency, memory, or other efficiency metrics. Multi-objective oracle could yield different rank selections and potentially larger gap. Adding efficiency constraints likely increases the oracle gap (more dimensions for optimization) but different optimal rank distribution. Our simplified metric provides a lower bound on the potential benefit.

\textbf{Limitation 4: Scope Limited to NLP on LLaMA-2-7B}

Results cover only 17 NLP tasks on a single decoder-only transformer (LLaMA-2-7B). Cross-modal generalization (vision, audio) and cross-architectural generalization (encoder-only, encoder-decoder) are unverified. Oracle gap pattern may not hold for other modalities or model families.

Within the NLP domain on decoder-only transformers, evidence is robust. 17 NLP tasks span sufficient diversity for existence proof: sentiment, NLI, paraphrase, similarity, cross-lingual transfer; dataset sizes from 635 to 392,702 samples; 9 linguistic phenomena. Cross-modal extension is future work, not a limitation of the core contribution.

\textbf{Limitation 5: Rank-32 Performance Reflects Uniform Protocol}

Rank-32 performs poorly (62.95\% average) but uses identical hyperparameters as other ranks. Rank-specific tuning (different learning rate, dropout, regularization) might improve rank-32 performance. We cannot definitively distinguish whether rank-32's poor performance reflects fundamental overfitting requiring >300K samples, or hyperparameter mismatch from uniform protocol.

Current evidence: Rank-32's collapse to 50\% on SST-2 (67K samples) and CoLA (8.5K samples) suggests fundamental overfitting rather than just hyperparameter issues. Literature consensus (Hu et al., 2021 uses different learning rates for different ranks; Zhang et al., 2023 adapts rank with rank-specific regularization) suggests rank-32 may require specialized tuning.

We acknowledge ambiguity: Our uniform protocol ensures fair comparison across ranks but creates this interpretative challenge. We cannot claim the protocol is "conservative" (underestimates gap) without evidence that tuning would improve rank-32.

\subsubsection{Broader Impact}

Improved resource efficiency in multi-domain deployment: Task-aware adapter configuration enables matching capacity to task complexity, avoiding over-parameterization for simple tasks and under-parameterization for complex tasks. Systems serving heterogeneous task distributions could reduce computational waste while improving average performance.

Informed adapter selection guidelines: Our finding that ranks 4-16 achieve consistent performance while rank-32 overfits on datasets below 300K samples provides empirical guidance for practitioners.

Foundation for adaptive configuration research: Quantifying the 15.09\% relative improvement oracle gap (upper bound) establishes a benchmark for future task-aware routing mechanisms. Researchers can now evaluate whether routing approaches justify added complexity by measuring what fraction of the gap they recover.

Increased system complexity: Task-aware routing adds components (task meta-feature extraction, routing classifier, adapter management) that increase deployment complexity. If routing introduces bugs or failures, it could degrade reliability below fixed-rank baselines despite higher oracle potential.

Routing errors could harm performance: Imperfect routing (selecting suboptimal rank) might perform worse than a safe fixed baseline. If routing accuracy <70\%, regret from wrong selections could dominate oracle gain, yielding negative net benefit.
